%\resizebox{\textwidth}{!}{%
\begin{tikzpicture}
\small

% Title
\node[dia/header=15cm] (title) {ІМЕРСИВНІ ТЕХНОЛОГІЇ В ОСВІТІ};

% Column headers — all at same y, minimum height ensures uniform height
\node[dia/rbox=3.0cm, minimum height=1.2cm, below=1.5cm of title.south, xshift=-5.7cm] (col1h) {\textbf{Імерсивні технології}};
\node[dia/rbox=3.0cm, minimum height=1.2cm, below=1.5cm of title.south, xshift=-1.9cm] (col2h) {\textbf{Імерсивні освітні ресурси}};
\node[dia/rbox=3.0cm, minimum height=1.2cm, below=1.5cm of title.south, xshift=1.9cm] (col3h) {\textbf{Цифровий дизайн в освіті}};
\node[dia/rbox=3.0cm, minimum height=1.2cm, below=1.5cm of title.south, xshift=5.7cm] (col4h) {\textbf{Педагогічна\\ ефективність ІОР}};

% Row a — all anchored to same y (relative to col1h bottom)
\node[dia/rbox=3.0cm, below=0.6cm of col1h] (col1a) {Віртуальна реальність};
\node[dia/rbox=3.0cm, anchor=center] at (col2h |- col1a) (col2a) {інтерактивність};
\node[dia/rbox=3.0cm, anchor=center] at (col3h |- col1a) (col3a) {інтерактивність};
\node[dia/rbox=3.0cm, anchor=center] at (col4h |- col1a) (col4a) {Засвоєння знань};

% Row b
\node[dia/rbox=3.0cm, below=0.6cm of col1a] (col1b) {Доповнена реальність};
\node[dia/rbox=3.0cm, anchor=center] at (col2h |- col1b) (col2b) {занурення};
\node[dia/rbox=3.0cm, anchor=center] at (col3h |- col1b) (col3b) {інтерактивність};
\node[dia/rbox=3.0cm, anchor=center] at (col4h |- col1b) (col4b) {Формування навичок};

% Row c
\node[dia/rbox=3.0cm, below=0.6cm of col1b] (col1c) {Змішана реальність};
\node[dia/rbox=3.0cm, anchor=center] at (col2h |- col1c) (col2c) {контекстуальність};
\node[dia/rbox=3.0cm, anchor=center] at (col3h |- col1c) (col3c) {інтерактивність};
\node[dia/rbox=3.0cm, anchor=center] at (col4h |- col1c) (col4c) {Мотивація до навчання};

% Vertical lines from title to column headers
\draw[dia/line] (title.south) -- ++(0,-0.5cm) -| (col1h.north);
\draw[dia/line] (title.south) -- ++(0,-0.5cm) -| (col2h.north);
\draw[dia/line] (title.south) -- ++(0,-0.5cm) -| (col3h.north);
\draw[dia/line] (title.south) -- ++(0,-0.5cm) -| (col4h.north);

% Horizontal arrows between sub-items
\draw[dia/arrow] (col1a.east) -- (col2a.west);
\draw[dia/arrow] (col2a.east) -- (col3a.west);
\draw[dia/arrow] (col3a.east) -- (col4a.west);

\draw[dia/arrow] (col1b.east) -- (col2b.west);
\draw[dia/arrow] (col2b.east) -- (col3b.west);
\draw[dia/arrow] (col3b.east) -- (col4b.west);

\draw[dia/arrow] (col1c.east) -- (col2c.west);
\draw[dia/arrow] (col2c.east) -- (col3c.west);
\draw[dia/arrow] (col3c.east) -- (col4c.west);

% Bottom section
\coordinate (botsep) at ($(col1c.south)!0.5!(col4c.south) + (0,-1.5cm)$);
\node[dia/header=15cm, anchor=north] at (botsep) (bottitle) {Процес навчання розробки ІОР};

% Three boxes in a chain
\node[dia/rbox=4cm, below=1.0cm of bottitle.south, xshift=-5.0cm] (step1)
  {Вивчення основ ІТ};
\node[dia/rbox=4cm, below=1.0cm of bottitle.south] (step2)
  {Освоєння інструментів розробки};
\node[dia/rbox=4cm, below=1.0cm of bottitle.south, xshift=5.0cm] (step3)
  {Педагогічне проєктування};

% Arrows between steps
\draw[dia/arrow] (step1.east) -- (step2.west);
\draw[dia/arrow] (step2.east) -- (step3.west);

% Lines from bottom title to steps
\draw[dia/line] (bottitle.south) -- ++(0,-0.3cm) -| (step1.north);
\draw[dia/line] (bottitle.south) -- ++(0,-0.3cm) -| (step2.north);
\draw[dia/line] (bottitle.south) -- ++(0,-0.3cm) -| (step3.north);

\end{tikzpicture}%
%}
